\documentclass[11pt,a4paper]{article}
\usepackage[margin=25mm]{geometry}
\usepackage{amsmath,amssymb,amsthm,bm,mathtools}
\usepackage[T1]{fontenc}
\usepackage{lmodern,microtype}
\usepackage[colorlinks=true,linkcolor=blue,urlcolor=blue]{hyperref}
\newcommand{\tr}{\operatorname{tr}}
\newcommand{\KL}{D_{\mathrm{KL}}}
\newcommand{\E}{\mathbb E}
\newcommand{\R}{\mathbb R}
\newcommand{\Id}{\mathrm I}
\newcommand{\N}{\mathcal N}
\newcommand{\F}{\mathcal F}
\newcommand{\Splus}{\mathbb S_{++}^{P}}
\newtheorem{proposition}{Proposition}
\title{Variational and information-theoretic interpretations\\of a latent-structured covariance equation}
\author{}
\date{}
\begin{document}
\maketitle
\begin{abstract}
The positive-definite solution of $CA^{-1}C-sC-B=0$ uniquely minimises
$\tr(A^{-1}C)+\tr(BC^{-1})-s\log\det C$.
We derive four complementary interpretations: an entropy--energy balance;
a compromise between Gaussian KL divergences in opposite directions;
reward for input-side latent overlap under covariance penalties; and,
after separating the isotropic background, genuine Gaussian-channel mutual-information maximisation under two costs.
We specify the random variables behind each information statement and distinguish information in a representation from information lost when the same matrix is interpreted as noise.
\end{abstract}

\section{Setup, domain, and two necessary qualifications}
Let
\begin{align}
 A&=S_X+aI_P,& S_X&=ZG_XZ^\top,& a&=\sigma_A^2,\\
 B&=S_Y+bI_P,& S_Y&=ZG_YZ^\top,& b&=\sigma_B^2,\\
 G_X&=W_X^\top K_XW_X,&
 G_Y&=W_Y^\top K_Y^{-1}W_Y.
\end{align}
Here $Z\in\R^{P\times K}$, $K<P$, the columns of $W_X,W_Y$ are orthonormal,
and $K_X,K_Y$ are positive definite. Consequently $G_X,G_Y\succ0$:
for example, $v^\top G_Xv=(W_Xv)^\top K_X(W_Xv)>0$ for $v\ne0$.
Throughout the main derivation, $a>0$, $b\ge0$, and $s>0$.

First, $Z$ cannot have full \emph{row} rank when $K<P$; we assume full
\emph{column} rank $K$. Second, the equation does not have a unique solution
over all symmetric matrices. Even in one dimension, for $A,B>0$, it has
one positive and one negative root. The relevant uniqueness statement is
\emph{uniqueness over $C\succ0$}, appropriate when $C$ is a covariance.
If $a=0$, the full-space inverse $A^{-1}$ does not exist. Section~\ref{sec:zero}
treats the reduced-space formulation and the limit $a\downarrow0$ separately.

All logarithms are natural. Additive constants in objectives are independent
of the optimisation variable. Covariances are expressed in fixed units;
a reference covariance can be inserted in every log determinant without
changing any minimiser.

\section{The basic convex variational principle}
\begin{proposition}
For $A\succ0$, $B\succeq0$, $s>0$, the functional
\begin{equation}\label{eq:F}
 \boxed{\F(C)=\tr(A^{-1}C)+\tr(BC^{-1})-s\log\det C,
 \qquad C\succ0}
\end{equation}
has a unique global minimiser, and its stationarity equation is
\begin{equation}\label{eq:riccati}
 CA^{-1}C-sC-B=0.
\end{equation}
\end{proposition}
\begin{proof}
For a symmetric variation $H$, use
$dC^{-1}[H]=-C^{-1}HC^{-1}$ and $d\log\det C[H]=\tr(C^{-1}H)$.
Then
\begin{equation}
 d\F(C)[H]=\tr\!\left[
 (A^{-1}-C^{-1}BC^{-1}-sC^{-1})H\right].
\end{equation}
Multiplying the vanishing gradient on the left and right by $C$ gives
\eqref{eq:riccati}. This does not require $A,B,C$ to commute.

For strict convexity, put $T=C^{-1/2}HC^{-1/2}$ and
$D=C^{-1/2}BC^{-1/2}\succeq0$. The second variation is
\begin{equation}
 d^2\F(C)[H,H]=2\tr(DT^2)+s\tr(T^2)>0\quad(H\ne0).
\end{equation}
The first term is nonnegative since both $D$ and $T^2$ are positive
semidefinite. Existence follows from coercivity: with
$m=\lambda_{\min}(A^{-1})>0$,
\begin{equation}
 \F(C)\ge\sum_{i=1}^P\bigl[m\lambda_i(C)-s\log\lambda_i(C)\bigr].
\end{equation}
Each scalar summand diverges at zero and infinity and is bounded below.
Thus sublevel sets stay away from the boundary and from infinity.
\end{proof}

For completeness, whiten by $A$:
\begin{equation}
 Q=A^{-1/2}CA^{-1/2},\qquad D=A^{-1/2}BA^{-1/2}.
\end{equation}
The equation becomes $Q^2-sQ-D=0$. Hence $Q$ commutes with $D$, and
positive definiteness selects the positive scalar root in every eigenspace:
\begin{equation}\label{eq:solution}
 \boxed{C^*=A^{1/2}\left[\frac{s}{2}I+
 \left(\frac{s^2}{4}I+A^{-1/2}BA^{-1/2}\right)^{1/2}\right]A^{1/2}.}
\end{equation}
For $D=0$ the positive root is $s$, not zero. In particular
$C^*\succeq sA$, with strict inequality if $B\succ0$.

\section{Interpretation I: maximum entropy with two quadratic costs}
Introduce independent auxiliary Gaussian vectors
$h\sim\N(0,C)$ and $y\sim\N(0,B)$ (the latter may be degenerate).
Then
\begin{align}
 \E[h^\top A^{-1}h]&=\tr(A^{-1}C),\\
 \E[y^\top C^{-1}y]&=\tr(BC^{-1}),\\
 \mathsf h(h)&=\tfrac12\log\det C+\tfrac P2\log(2\pi e).
\end{align}
Consequently
\begin{equation}\label{eq:entropy}
 \F(C)=\E\|h\|_{A^{-1}}^2+\E\|y\|_{C^{-1}}^2
 -2s\mathsf h(h)+\text{constant}.
\end{equation}
This is a maximum-entropy covariance subject to two competing costs.
The first controls representation variance in the metric supplied by $A$:
variance in a large-eigenvalue direction of $A$ is relatively cheap.
The second prevents $C$ from assigning too little variance to directions
that occur in $B$. Entropy rewards a large occupied volume and prevents
collapse, even in directions where $B$ vanishes.

Equivalently, let $e_X^*=\tr(A^{-1}C^*)$ and
$e_Y^*=\tr(B(C^*)^{-1})$. Then $C^*$ uniquely maximises
$\log\det C$ subject to
\begin{equation}
 \tr(A^{-1}C)\le e_X^*,\qquad \tr(BC^{-1})\le e_Y^*.
\end{equation}
Indeed, a feasible covariance with a larger log determinant would have
smaller $\F$, contradicting optimality. The parameter $s$ sets the
entropy reward in the penalised formulation; it does not by itself specify
either of these two budgets independently.

The latent structure enters both costs. Writing
$y=ZG_Y^{1/2}v+\sqrt b\,\epsilon$ with independent standard normal vectors,
\begin{equation}\label{eq:Ytrace}
 \tr(BC^{-1})=
 \|C^{-1/2}ZG_Y^{1/2}\|_F^2+b\tr(C^{-1}).
\end{equation}
Thus the output-side cost measures how large the latent signal remains
after whitening by the candidate covariance. Minimising it encourages
coverage of that signal by $C$. It is not, without an additional observation
model, a latent reconstruction error.

\section{Interpretation II: Gaussian KL compromise}
For positive-definite $U,V$,
\begin{equation}\label{eq:KL}
 2\KL\bigl(\N(0,U)\Vert\N(0,V)\bigr)
 =\tr(V^{-1}U)-P+\log\det V-\log\det U.
\end{equation}
Taking $U=C$, $V=sA$ gives the exact identity, up to a constant,
\begin{equation}\label{eq:oneKL}
 \boxed{\F(C)=2s\KL\bigl(\N(0,C)\Vert\N(0,sA)\bigr)
 +\tr(BC^{-1})+\text{constant}.}
\end{equation}
So $C$ stays close to an input-side reference covariance $sA$, while
expanding to cover output-side fluctuations. This form remains valid for
$b=0$ because it never puts a degenerate $B$ into a Gaussian KL divergence.

If $B\succ0$, choose any $r>0$ and set $t=r+s$. Then
\begin{align}
 2t\KL\bigl(\N(0,C)\Vert\N(0,tA)\bigr)
 &=\tr(A^{-1}C)-t\log\det C+\text{constant},\\
 2r\KL\bigl(\N(0,B/r)\Vert\N(0,C)\bigr)
 &=\tr(BC^{-1})+r\log\det C+\text{constant}.
\end{align}
Adding yields
\begin{equation}\label{eq:twoKL}
 \boxed{\F(C)=2t\KL\bigl(\N(0,C)\Vert\N(0,tA)\bigr)
 +2r\KL\bigl(\N(0,B/r)\Vert\N(0,C)\bigr)+\text{constant}.}
\end{equation}
For example $r=s$, $t=2s$ produces weights $4s,2s$ and references
$2sA,B/s$. The rescaling of the reference covariances is essential:
an arbitrary weighted sum using unscaled $A,B$ does not generally give
the original equation.

The first divergence penalises putting representation variance where the
input reference has little variance. The second penalises failing to cover
variance present in the output reference. Both refer to the same latent
subspace but use different latent metrics $G_X,G_Y$.
The freedom in $r$ shows that this interpretation is not a uniquely
identified probabilistic model; these are equivalent objective decompositions.
For $b=0$ the full-space second KL is infinite, so use \eqref{eq:oneKL}
or the reduced-space version below.

\section{Interpretation III: reward for input-side latent overlap}
This formulation makes explicit the reward for alignment with $Z$.
Define the symmetric matrix
\begin{equation}
 H_X=S_X(S_X+aI)^{-1}=I-aA^{-1}.
\end{equation}
It is supported on $\operatorname{col}(Z)$, has eigenvalues in $[0,1)$,
and assigns larger weight to stronger input-side signal directions.
Since $A^{-1}=a^{-1}(I-H_X)$, equation~\eqref{eq:F} becomes
\begin{equation}\label{eq:overlap}
 \boxed{\F(C)=\frac1a\tr C-\frac1a\tr(H_XC)
 +\tr(BC^{-1})-s\log\det C.}
\end{equation}
Equivalently, $C^*$ maximises
\begin{equation}
 \underbrace{\tr(H_XC)}_{\text{input latent overlap}}
 +\underbrace{as\log\det C}_{\text{volume reward}}
 -\underbrace{\tr C}_{\text{total variance cost}}
 -\underbrace{a\tr(BC^{-1})}_{\text{output coverage cost}}.
\end{equation}
This is a direct overlap-maximisation statement, although the exact overlap
is with a noise-weighted latent kernel, not generally with $ZZ^\top$.

To see its latent-coordinate form, let
\begin{equation}
 V_X=\left(G_X^{-1}+a^{-1}Z^\top Z\right)^{-1}.
\end{equation}
Woodbury gives
\begin{equation}
 A^{-1}=a^{-1}I-a^{-2}ZV_XZ^\top,
 \qquad \tr(H_XC)=a^{-1}\tr(V_XZ^\top CZ).
\end{equation}
Thus $Z^\top CZ$ is rewarded through a specific latent metric.
Only when that metric is proportional to the identity is this proportional
to the unweighted overlap $\tr(Z^\top CZ)$.

There is also a posterior interpretation. In the input-side model
\begin{equation}
 u\sim\N(0,G_X),\qquad h\mid u\sim\N(Zu,aI),
\end{equation}
the posterior is $u\mid h\sim\N(m(h),V_X)$, with
$m(h)=a^{-1}V_XZ^\top h$. Averaging over the candidate law $h\sim\N(0,C)$,
\begin{equation}
 \E_C[m(h)^\top V_X^{-1}m(h)]
 =a^{-2}\tr(V_XZ^\top CZ)=a^{-1}\tr(H_XC).
\end{equation}
The overlap reward is therefore the mean squared inferred latent amplitude,
measured in posterior-uncertainty units. This is posterior activation,
not automatically accuracy relative to a true latent variable: the averaging
law is $\N(0,C)$, which need not equal the input model's marginal $\N(0,A)$.

Equivalently, the log Bayes factor of the structured input model against
pure isotropic noise is
\begin{equation}
 \log\frac{p_A(h)}{p_{aI}(h)}
 =\frac1{2a}h^\top H_Xh-\frac12\log\det(A/a).
\end{equation}
Hence the negative overlap term in \eqref{eq:overlap} is exactly
$-2\E_C\log[p_A(h)/p_{aI}(h)]$ plus a constant. It rewards samples whose
fluctuations provide evidence for the input latent model.

\section{The exact latent-subspace reduction}\label{sec:reduction}
Put $R=(Z^\top Z)^{1/2}$ and $U=ZR^{-1}$, so that $U^\top U=I_K$,
$Z=UR$, and $\Pi=UU^\top$ projects onto the latent subspace.
Define
\begin{equation}
 A_L=RG_XR+aI_K,\qquad B_L=RG_YR+bI_K.
\end{equation}
In the decomposition $\R^P=\operatorname{col}(Z)\oplus\operatorname{col}(Z)^\perp$,
$A=A_L\oplus aI$ and $B=B_L\oplus bI$.
Formula~\eqref{eq:solution} therefore implies
\begin{equation}\label{eq:blocks}
 C^*=UC_L^*U^\top+c_0(I-\Pi),\qquad
 c_0=\frac{sa+\sqrt{s^2a^2+4ab}}2,
\end{equation}
where $C_L^*$ uniquely minimises
\begin{equation}\label{eq:FL}
 \F_L(C_L)=\tr(A_L^{-1}C_L)+\tr(B_LC_L^{-1})-s\log\det C_L.
\end{equation}
The complementary contribution is
$(P-K)[c_0/a+b/c_0-s\log c_0]$.
All non-isotropic structure lies in the $K$-dimensional latent subspace,
irrespective of the ambient dimension $P$.

An additional useful fact is
\begin{equation}\label{eq:floor}
 C_L^*\succ c_0I_K.
\end{equation}
To prove it, take a unit eigenvector $v$ of $C_L^*$ with eigenvalue $c>0$.
Since $A_L\succ aI$ and $B_L\succ bI$,
\begin{equation}
 c^2v^\top A_L^{-1}v=sc+v^\top B_Lv,
 \qquad \frac{c^2}{a}>sc+b.
\end{equation}
The positive root of $c^2/a-sc-b=0$ is $c_0$, so $c>c_0$.
Consequently there is a unique $M^*\succ0$ such that
\begin{equation}\label{eq:M}
 C^*=c_0I+ZM^*Z^\top,\qquad
 M^*=R^{-1}(C_L^*-c_0I)R^{-1}.
\end{equation}
The value of $c_0$ is determined explicitly from $a,b,s$; finding it does
not require knowing $C^*$.

\section{Interpretation IV: information maximisation in a latent channel}\label{sec:MI}
The preceding reduction supplies a precise positive mutual-information
interpretation. Restrict to
\begin{equation}
 C(M)=c_0I+ZMZ^\top,\qquad M\succeq0.
\end{equation}
This restriction preserves the global optimum by \eqref{eq:M}.
Introduce a fixed standard Gaussian latent vector $u\in\R^K$ and a
representation $h\in\R^P$ through
\begin{equation}\label{eq:channel}
 u\sim\N(0,I_K),\qquad
 h=ZM^{1/2}u+\epsilon,\qquad
 \epsilon\sim\N(0,c_0I_P),\quad \epsilon\perp u.
\end{equation}
Its marginal covariance is $C(M)$ and its conditional noise covariance
does not depend on $M$. Subtracting the two Gaussian entropies gives
\begin{align}
 \Id_M(u;h)
 &=\mathsf h(h)-\mathsf h(h\mid u)\\
 &=\frac12\log\det C(M)-\frac P2\log c_0\\
 &=\frac12\log\det\left(I_K+c_0^{-1}M^{1/2}Z^\top ZM^{1/2}\right).
\end{align}
Thus an exact variational principle for the same $C^*$ is
\begin{equation}\label{eq:MIobjective}
 \boxed{M^*=\arg\min_{M\succeq0}\left\{
 \tr(A^{-1}C(M))+\tr(BC(M)^{-1})-2s\Id_M(u;h)\right\}.}
\end{equation}
This is not a bound or a large-noise approximation.
Although the parametrisation of the channel uses a square root, the objective
is strictly convex in $M$: it equals $\F(C(M))$ plus a constant, and
$M\mapsto ZMZ^\top$ is injective.

\paragraph{What is being learned?}
The optimisation chooses the strength and mixing of latent signal transmitted
through the fixed patterns in $Z$. Mutual information rewards a representation
from which the Gaussian coefficients $u$ can be inferred. The first trace
charges for representation variance according to the input-side metric.
The second discourages allocating too little covariance to output-relevant
latent directions. The tradeoff therefore selects an informative latent
representation that respects both input and output structure.

This also has an exact uncertainty interpretation. Completing the square gives
\begin{equation}
 \operatorname{Cov}(u\mid h)
 =\left(I_K+c_0^{-1}M^{1/2}Z^\top ZM^{1/2}\right)^{-1},
\end{equation}
so $2\Id_M(u;h)=-\log\det\operatorname{Cov}(u\mid h)$.
The reward reduces the posterior uncertainty volume of the latent
coefficients. It need not maximise every coordinate's accuracy independently.

As in Section 3, at the attained budgets $e_X^*,e_Y^*$, $M^*$ also solves
\begin{equation}
 \max_{M\succeq0}\ \Id_M(u;h)
 \quad\text{subject to}\quad
 \tr(A^{-1}C(M))\le e_X^*,\quad
 \tr(BC(M)^{-1})\le e_Y^*.
\end{equation}
Information alone would increase without bound as signal strength grows;
the trace costs are an essential part of this interpretation.

\paragraph{Scope of the information statement.}
Here $Z$ is a fixed loading matrix, not a random variable. The mutual
information is between random latent coefficients and their representation
\emph{conditional on that fixed design}. It is not $\Id(Z;C)$, nor is it
automatically mutual information between a randomly sampled row $z^\mu$
and an observed feature. Those quantities require a prior over $Z$, or a
specified sampling and observation model for its rows.
The original equation alone supplies neither. The channel in
\eqref{eq:channel} is an explicit, compatible interpretation of the objective,
not a uniquely implied data-generating process.

\section{Why a different channel can reverse the information interpretation}
To avoid confusing covariance-as-signal with covariance-as-noise, consider
instead a fixed output-side latent signal with \emph{noise} covariance $C$:
\begin{equation}
 v\sim\N(0,I_K),\qquad
 \widetilde y=ZG_Y^{1/2}v+\eta_C,\qquad \eta_C\sim\N(0,C).
\end{equation}
Set $T=C^{-1/2}S_YC^{-1/2}$. Direct Gaussian entropy and KL calculations give
\begin{align}
 2\Id(v;\widetilde y)&=\log\det(I+T),\\
 2\KL\bigl(\N(0,C+S_Y)\Vert\N(0,C)\bigr)
 &=\tr T-\log\det(I+T).
\end{align}
Hence the identity
\begin{equation}
 \tr(S_YC^{-1})=2\Id(v;\widetilde y)
 +2\KL\bigl(\N(0,C+S_Y)\Vert\N(0,C)\bigr)
\end{equation}
is exact, and
\begin{align}
 \F(C)={}&2s\KL\bigl(\N(0,C)\Vert\N(0,sA)\bigr)
 +2\Id(v;\widetilde y)\nonumber\\
 &+2\KL\bigl(\N(0,C+S_Y)\Vert\N(0,C)\bigr)
 +b\tr(C^{-1})+\text{constant}.
\end{align}
Here the information term has a positive sign in a minimisation objective:
larger $C$ masks a fixed signal. There is no contradiction with
Section~\ref{sec:MI}, where $C$ is the \emph{total representation covariance}
and the observation noise is fixed at $c_0I$. One must specify which of
these roles $C$ plays before claiming information preservation or suppression.

\section{Modes, invariances, and the limits of latent identification}
If $A_L$ and $B_L$ commute, let their simultaneous eigenvalues be
$\alpha_i>0,\beta_i>0$. Then
\begin{equation}
 c_i^*=\frac{s\alpha_i+
 \sqrt{s^2\alpha_i^2+4\alpha_i\beta_i}}2,
 \qquad
 \F_L=\sum_i\left[\frac{c_i}{\alpha_i}+\frac{\beta_i}{c_i}-s\log c_i\right].
\end{equation}
Input variance makes a mode cheap, output variance makes its coverage
valuable, and entropy opposes collapse. Each $c_i^*$ increases with both
$\alpha_i$ and $\beta_i$. In this case the channel information at the optimum is
\begin{equation}
 \Id(u;h)=\frac12\sum_{i=1}^K\log(c_i^*/c_0).
\end{equation}
The determinant formula remains valid without commutation, but there need
not be a common Euclidean eigenbasis for the input and output kernels.

For $a>0$, full-column-rank $Z$ gives
$\operatorname{range}(C^*-c_0I)=\operatorname{col}(Z)$.
Thus the latent subspace is recoverable from the exact optimum, provided
the background level is known. The coordinates of $Z$ are not identifiable:
for any invertible $T\in\R^{K\times K}$, the simultaneous transformations
\begin{equation}
 Z\mapsto ZT,\qquad
 G_X\mapsto T^{-1}G_XT^{-\top},\qquad
 G_Y\mapsto T^{-1}G_YT^{-\top}
\end{equation}
leave $A,B,C^*$ unchanged. Consequently the equation encodes latent
geometry through $S_X,S_Y$, not a distinguished set of latent coordinates.

\section{Zero-noise cases}\label{sec:zero}
\paragraph{$a>0$, $b=0$.}
The original equation and convex functional remain well defined.
Here $B$ is singular, but strict convexity and existence still hold because
$s>0$. The complementary variance is $c_0=sa>0$, so the exact channel
interpretation remains valid. Use \eqref{eq:oneKL}, rather than a full-space
KL with singular $B$ as its first covariance. On the latent subspace,
$B_L=RG_YR\succ0$, so the two-KL formulation is valid there.

\paragraph{$a=0$, with $b\ge0$.}
The original full-space equation is undefined, since $A=S_X$ has rank $K<P$.
One well-defined replacement is the equation on the latent subspace:
\begin{equation}
 C_L(RG_XR)^{-1}C_L-sC_L-(RG_YR+bI_K)=0.
\end{equation}
Its strictly convex functional is \eqref{eq:FL} with $a=0$, and all
positive-definite Gaussian and entropy arguments apply in dimension $K$.

Alternatively, take the solutions with $a>0$ and then let $a\downarrow0$.
The latent block converges to the solution just displayed, while
$c_0\to0$. The full-space limit is singular and supported on
$\operatorname{col}(Z)$; it is not a solution involving the nonexistent
ordinary inverse of $S_X$. Objective values on the complementary subspace
can diverge and must be removed as constants before comparing reduced
objectives. Merely replacing $A^{-1}$ by $A^+$ is a different problem:
for a supported candidate $C$, its complementary equation would read
$-bI=0$, which fails when $b>0$.

The specific mutual-information model of Section~\ref{sec:MI} uses
$c_0>0$. As $a\downarrow0$, its noise vanishes and its mutual information
diverges for a nondegenerate continuous latent signal. Thus it does not
give a finite information value at exactly zero noise. A reduced-space
channel with an explicitly chosen positive noise floor can be introduced,
but that is an additional modelling choice rather than a finite noiseless
limit of the displayed channel.

\section{Which formulation best answers the latent-information question?}
The common mathematical core is the strictly convex objective
\eqref{eq:F}. For covariance fitting, \eqref{eq:oneKL} is particularly
direct: retain input-side Gaussian structure while covering output-side
variance. For an input-centred geometric interpretation,
\eqref{eq:overlap} explicitly rewards the latent overlap
$\tr(V_XZ^\top CZ)$, balanced against total variance, output coverage,
and entropy. For a literal positive information objective,
\eqref{eq:MIobjective} is exact: it maximises information transmitted
through the latent patterns in $Z$ at a fixed isotropic noise level,
subject to the two covariance costs.

These interpretations are compatible, but they answer different questions.
None licenses an unqualified claim that $C$ maximises information about
an unspecified random $Z$. Once the latent observation model is stated,
the precise information quantity, its sign, and its regularising costs
follow explicitly.
\end{document}
